\documentclass[11pt, a4paper]{article}

% ── Packages ──────────────────────────────────────────────────────────
\usepackage[utf8]{inputenc}
\usepackage[T1]{fontenc}
\usepackage{lmodern}
\usepackage[margin=1in]{geometry}
\usepackage{amsmath, amssymb, amsthm}
\usepackage{booktabs}
\usepackage{longtable}
\usepackage{graphicx}
\usepackage{hyperref}
\usepackage{xcolor}
\usepackage{listings}
\usepackage{tcolorbox}
\usepackage{polya}

\hypersetup{
  colorlinks=true,
  linkcolor=polyablue,
  urlcolor=polyablue,
  citecolor=polyablue,
}

% ── Code listing style ────────────────────────────────────────────────
\lstdefinestyle{polya-python}{
  language=Python,
  basicstyle=\ttfamily\small,
  keywordstyle=\color{polyablue}\bfseries,
  commentstyle=\color{polyagray}\itshape,
  stringstyle=\color{polyagreen},
  numbers=left,
  numberstyle=\tiny\color{polyagray},
  numbersep=8pt,
  frame=single,
  framerule=0.4pt,
  rulecolor=\color{polyagray},
  breaklines=true,
  showstringspaces=false,
  tabsize=4,
}
\lstset{style=polya-python}
\title{Nash Equilibrium -- Coffee Shop Pricing}
\author{uber-polya}
\date{2026-02-26}

\begin{document}

\maketitle
\tableofcontents
\newpage

% ── Phase A: Problem Formulation ─────────────────────────────────────
\section{Problem Formulation}

\begin{polyamodel}

\textbf{Problem Type}: Find \hfill
\textbf{Domain}: Game Theory

\textbf{Named Problem}: Game of Chicken

\end{polyamodel}

% ── Universe ──────────────────────────────────────────────────────────
\subsection{Universe}

\begin{itemize}
  \item $Strategies: \{Discount, Hold\}$
  \item $Strategies: \{Discount, Hold, Loyalty\}$
\end{itemize}

% ── Variables ─────────────────────────────────────────────────────────
\subsection{Variables}

\begin{longtable}{lll}
\toprule
\textbf{Symbol} & \textbf{Meaning} & \textbf{Type / Range} \\
\midrule
\endhead
$\sigma_i$ & mixed strategy for player i (probability vector) & $\Delta(S_i)$ \\
$A, B$ & payoff matrices for players 1 and 2 & $\mathbb{R}^{m \times n}$ \\
$u_i(\sigma)$ & expected utility of player i under profile sigma & $\mathbb{R}$ \\
\bottomrule
\end{longtable}

% ── Structure ─────────────────────────────────────────────────────────
\subsection{Structure}

Two competing coffee shops on the same street decide simultaneously whether to run a promotion (discount) or hold prices. Each shop's profit depends on both its own action and its competitor's action:

- If both discount, they trigger a destructive price war (\$2k each).
- If one discounts while the other holds, the discounter captures market share (\$5k) while the holder loses traffic (\$3k).
- If both hold prices, they share the market comfortably (\$4k each).

This is an anti-coordination game (Game of Chicken variant): each shop wants to do the opposite of its competitor, but they move simultaneously without knowing the other's choice.

A second, richer variant adds a third strategy -- "loyalty program" -- creating a 3x3 symmetric game with circular dominance (Discount beats Hold, Hold beats Loyalty, Loyalty beats Discount), analogous to Rock-Paper-Scissors but with cooperative payoffs on the diagonal.

% ── Mapping ───────────────────────────────────────────────────────────
\subsection{Real-World Mapping}

\begin{longtable}{ll}
\toprule
\textbf{Real-World Concept} & \textbf{Mathematical Object} \\
\midrule
\endhead
coffee shop & $player i$ \\
pricing decision & $strategy s_i in S_i$ \\
profit & $payoff u_i(s)$ \\
payoff table & $bimatrix (A, B)$ \\
\bottomrule
\end{longtable}

% ── Constraints ───────────────────────────────────────────────────────
\subsection{Constraints}

\begin{enumerate}
  \item $\sigma_i \ge 0, \quad \sum_{k} \sigma_{ik} = 1$
  \hfill \textit{(mixed strategy is a probability distribution)}
  \item $u_i(\sigma_i^*, \sigma_{-i}^*) \ge u_i(s_i, \sigma_{-i}^*) \quad \forall s_i \in S_i$
  \hfill \textit{(Nash equilibrium: no profitable unilateral deviation)}
  \item $A \in \mathbb{R}^{m \times n}, \quad B \in \mathbb{R}^{m \times n}$
  \hfill \textit{(payoff matrices define the normal-form game)}
\end{enumerate}

% ── Objective / Claim ─────────────────────────────────────────────────
\subsection{Objective}

\begin{equation}
\text{Find } \sigma^* \text{ s.t. } u_i(\sigma_i^*, \sigma_{-i}^*) \ge u_i(s_i, \sigma_{-i}^*) \; \forall i, s_i
\end{equation}

% ── Approach ──────────────────────────────────────────────────────────
\subsection{Approach}

\begin{description}
  \item[Suggested approach:] unknown
  \item[Complexity class:] 
  \item[Available tools:] unknown
\end{description}
% ── Phase B: Solution ────────────────────────────────────────────────
\section{Solution}

\begin{polyaresult}

\textbf{Answer}: 2x2 Coffee Shop Pricing (Anti-coordination): 3 Nash equilibria found; 3x3 Coffee Shop Pricing (Circular Dominance with Loyalty): 1 Nash equilibria found


\textbf{Optimal}: No / Unknown \hfill
\textbf{Feasible}: Yes
\textbf{Algorithm}: unknown \quad $$

\textbf{Time}: 0.0s

\textbf{Certificate}: Nash equilibria verified: no player has a profitable unilateral deviation.

\end{polyaresult}

% ── Solution Details ──────────────────────────────────────────────────
\subsection{Solution Details}

Game: 2x2 Coffee Shop Pricing (Anti-coordination)

Payoff matrix (Player 1, Player 2):
        Discount           Hold
    Discount      (2.0, 2.0)     (5.0, 3.0)
        Hold      (3.0, 5.0)     (4.0, 4.0)

Nash equilibria found: 3
  NE 1 (Pure): sigma\_1 = [1.0, 0.0], sigma\_2 = [0.0, 1.0]
    Payoffs: (5.0, 3.0)
    Support: P1 plays Discount; P2 plays Hold
  NE 2 (Pure): sigma\_1 = [0.0, 1.0], sigma\_2 = [1.0, 0.0]
    Payoffs: (3.0, 5.0)
    Support: P1 plays Hold; P2 plays Discount
  NE 3 (Mixed): sigma\_1 = [0.5, 0.5], sigma\_2 = [0.5, 0.5]
    Payoffs: (3.5, 3.5)
    Support: P1 plays Discount, Hold; P2 plays Discount, Hold
  No dominant strategy for Shop 1
  No dominant strategy for Shop 2
  Zero-sum: No
  Minimax value (Player 1): 3.0
  Minimax value (Player 2): 3.0

Game: 3x3 Coffee Shop Pricing (Circular Dominance with Loyalty)

Payoff matrix (Player 1, Player 2):
        Discount           Hold        Loyalty
    Discount      (3.0, 3.0)     (6.0, 1.0)     (2.0, 5.0)
        Hold      (1.0, 6.0)     (5.0, 5.0)     (7.0, 2.0)
     Loyalty      (5.0, 2.0)     (2.0, 7.0)     (4.0, 4.0)

Nash equilibria found: 1
  NE 1 (Mixed): sigma\_1 = [0.42857142857142855, 0.3333333333333333, 0.2380952380952381], sigma\_2 = [0.42857142857142855, 0.3333333333333333, 0.2380952380952381]
    Payoffs: (3.7619047619047614, 3.761904761904762)
    Support: P1 plays Discount, Hold, Loyalty; P2 plays Discount, Hold, Loyalty
  No dominant strategy for Shop 1
  No dominant strategy for Shop 2
  Zero-sum: No
  Minimax value (Player 1): 3.7619047619047614
  Minimax value (Player 2): 3.7619047619047614


% ── Verification ──────────────────────────────────────────────────────
\subsection{Verification}

\begin{longtable}{lcl}
\toprule
\textbf{Check} & \textbf{Status} & \textbf{Detail} \\
\midrule
\endhead
game\_1\_2x2: eq0 probs sum 1 p1 & \passmark & Passed \\
game\_1\_2x2: eq0 probs sum 1 p2 & \passmark & Passed \\
game\_1\_2x2: eq0 nonneg p1 & \passmark & Passed \\
game\_1\_2x2: eq0 nonneg p2 & \passmark & Passed \\
game\_1\_2x2: eq0 payoff p1 correct & \passmark & Passed \\
game\_1\_2x2: eq0 payoff p2 correct & \passmark & Passed \\
game\_1\_2x2: eq0 no profitable deviation p1 & \passmark & Passed \\
game\_1\_2x2: eq0 no profitable deviation p2 & \passmark & Passed \\
game\_1\_2x2: eq0 indifference condition p1 & \passmark & Passed \\
game\_1\_2x2: eq0 indifference condition p2 & \passmark & Passed \\
game\_1\_2x2: eq1 probs sum 1 p1 & \passmark & Passed \\
game\_1\_2x2: eq1 probs sum 1 p2 & \passmark & Passed \\
game\_1\_2x2: eq1 nonneg p1 & \passmark & Passed \\
game\_1\_2x2: eq1 nonneg p2 & \passmark & Passed \\
game\_1\_2x2: eq1 payoff p1 correct & \passmark & Passed \\
game\_1\_2x2: eq1 payoff p2 correct & \passmark & Passed \\
game\_1\_2x2: eq1 no profitable deviation p1 & \passmark & Passed \\
game\_1\_2x2: eq1 no profitable deviation p2 & \passmark & Passed \\
game\_1\_2x2: eq1 indifference condition p1 & \passmark & Passed \\
game\_1\_2x2: eq1 indifference condition p2 & \passmark & Passed \\
game\_1\_2x2: eq2 probs sum 1 p1 & \passmark & Passed \\
game\_1\_2x2: eq2 probs sum 1 p2 & \passmark & Passed \\
game\_1\_2x2: eq2 nonneg p1 & \passmark & Passed \\
game\_1\_2x2: eq2 nonneg p2 & \passmark & Passed \\
game\_1\_2x2: eq2 payoff p1 correct & \passmark & Passed \\
game\_1\_2x2: eq2 payoff p2 correct & \passmark & Passed \\
game\_1\_2x2: eq2 no profitable deviation p1 & \passmark & Passed \\
game\_1\_2x2: eq2 no profitable deviation p2 & \passmark & Passed \\
game\_1\_2x2: eq2 indifference condition p1 & \passmark & Passed \\
game\_1\_2x2: eq2 indifference condition p2 & \passmark & Passed \\
game\_1\_2x2: num equilibria positive & \passmark & Passed \\
game\_1\_2x2: support enum found all & \passmark & Passed \\
game\_1\_2x2: zero sum correct & \passmark & Passed \\
game\_1\_2x2: all checks pass & \passmark & Passed \\
game\_2\_3x3: eq0 probs sum 1 p1 & \passmark & Passed \\
game\_2\_3x3: eq0 probs sum 1 p2 & \passmark & Passed \\
game\_2\_3x3: eq0 nonneg p1 & \passmark & Passed \\
game\_2\_3x3: eq0 nonneg p2 & \passmark & Passed \\
game\_2\_3x3: eq0 payoff p1 correct & \passmark & Passed \\
game\_2\_3x3: eq0 payoff p2 correct & \passmark & Passed \\
game\_2\_3x3: eq0 no profitable deviation p1 & \passmark & Passed \\
game\_2\_3x3: eq0 no profitable deviation p2 & \passmark & Passed \\
game\_2\_3x3: eq0 indifference condition p1 & \passmark & Passed \\
game\_2\_3x3: eq0 indifference condition p2 & \passmark & Passed \\
game\_2\_3x3: num equilibria positive & \passmark & Passed \\
game\_2\_3x3: support enum found all & \passmark & Passed \\
game\_2\_3x3: zero sum correct & \passmark & Passed \\
game\_2\_3x3: all checks pass & \passmark & Passed \\
\bottomrule
\end{longtable}

% ── Phase C: Interpretation ──────────────────────────────────────────
\section{Interpretation}

% ── The Question & The Answer ─────────────────────────────────────────
\subsection{The Question}
Two competing coffee shops on the same street decide simultaneously whether to run a promotion (discount) or hold prices.

\subsection{The Answer}

\begin{polyainsight}[title={Bottom Line}]
2x2 Coffee Shop Pricing (Anti-coordination): 2 pure + 1 mixed NE; 3x3 Coffee Shop Pricing (Circular Dominance with Loyalty): 0 pure + 1 mixed NE
\end{polyainsight}

% ── What This Means ───────────────────────────────────────────────────
\subsection{What This Means}

2x2 Game (Discount vs Hold -- Anti-coordination):
- No strictly dominant strategy for either player
- Two pure-strategy Nash equilibria: (Discount, Hold) at (\$5k, \$3k) and (Hold, Discount) at (\$3k, \$5k)
- One mixed-strategy Nash equilibrium: each shop plays 50/50, expected payoff \$3.50k each
- Security level: \$3.00k/month (guaranteed by pure Hold maximin strategy)

3x3 Game (Discount vs Hold vs Loyalty -- Circular Dominance):
- No strictly dominant or weakly dominated strategy
- Circular best-response structure: Loyalty beats Discount beats Hold beats Loyalty
- No pure-strategy Nash equilibrium (every pure profile is exploitable)
- Unique fully-mixed NE: (3/7 Discount, 1/3 Hold, 5/21 Loyalty)
- Expected payoff at NE: 79/21 \textasciitilde{} \$3.762k/month each
- Pareto-dominated by several pure outcomes (e.g., Hold-Hold at \$5k each)

1. Formulate as bimatrix game (A, B) where A[i,j] is Player 1's payoff and B[i,j] is Player 2's payoff when Player 1 plays row i and Player 2 plays column j
2. Enumerate Nash equilibria via support enumeration (nashpy): for each pair of support sets, solve the indifference conditions
3. Detect strictly dominant and weakly dominated strategies
4. Compute best response correspondences for each pure strategy
5. Compute maximin (security level) strategies via linear programming (scipy)
6. Analyze Pareto efficiency of equilibrium outcomes
7. Verify independently that no player has a profitable unilateral deviation
8. Cross-validate with vertex enumeration to confirm all equilibria found

% ── Sensitivity Analysis ──────────────────────────────────────────────

% ── Figures ───────────────────────────────────────────────────────────

% ── Recommendations ───────────────────────────────────────────────────
\subsection{Recommendations}

\begin{enumerate}
  \item Consider pre-play communication or binding agreements to reach Pareto-superior outcomes.
\end{enumerate}

% ── Limitations ───────────────────────────────────────────────────────
\subsection{Limitations}

\begin{polyawarnbox}
\begin{itemize}
  \item Assumes rational, utility-maximizing players with complete information.
  \item Real-world players may use heuristics or have incomplete information.
\end{itemize}
\end{polyawarnbox}

% ── Appendix ─────────────────────────────────────────────────────────

\end{document}